% Gemeinsame Originalaussagen; nur Band 46 registriert sie.
\newcommand{\FranklStatement}[1]{\csname FranklStatement@#1\endcsname}

% Rein typographische Anpassung der langen Definitionsformel im Fachband.
% FormulaDefDeltaKR registriert weiterhin die unveränderte Originalformel.
\newsavebox{\FranklCanonicalFormulaBox}
\newcommand{\FranklCanonicalFittedFormula}[1]{%
  \sbox{\FranklCanonicalFormulaBox}{\(\displaystyle #1\)}%
  \ifdim\wd\FranklCanonicalFormulaBox>\linewidth
    \resizebox{\linewidth}{!}{\usebox{\FranklCanonicalFormulaBox}}%
  \else
    \usebox{\FranklCanonicalFormulaBox}%
  \fi
}
\newcommand{\FranklCanonicalDefinitionDisplay}{%
  \RenewDocumentCommand{\FormulaDefDeltaK}{o m m +m}{%
    \IfNoValueTF{##1}{%
      \FormulaDefDeltaKR{\FranklCanonicalFittedFormula{##2}}{##2}{##3}{##4}%
    }{%
      \FormulaDefDeltaKR[##1]{\FranklCanonicalFittedFormula{##2}}{##2}{##3}{##4}%
    }%
  }%
}

\expandafter\newcommand\csname FranklStatement@FranklSingletonMember\endcsname{%
\FormulaThmDelta[Frankl-Vermutung bei enthaltener Einermenge]{%
  \Finite{M}\dsep \UnionClosedFam(M)\dsep \{a\}\in M\vdash \Frankl{M}
}{%
  \DeltaRow{Mengen}{M\dsep a}
  \DeltaRow{Funktionen}{\FranklAdj{M}{a}}
}
}

\expandafter\newcommand\csname FranklStatement@FranklAvoidingEmpty\endcsname{%
\FormulaThmDelta[Leere vermeidende Teilfamilie liefert Frankl]{%
  a\in\bigcup M\dsep \FamAvoiding{M}{a}=\varnothing
  \vdash \Frankl{M}
}{%
  \DeltaRow{Mengen}{M\dsep a}
  \DeltaRow{Funktionen}{F}
  \DeltaRow{Teilmengenfamilien}{%
    \FamAvoiding{M}{a}\dsep \FamContaining{M}{a}}
}
}

\expandafter\newcommand\csname FranklStatement@FranklNonemptyIntersection\endcsname{%
\FormulaThmDelta[Frankl-Vermutung bei nichtleerem Durchschnitt]{%
  \FranklFam{M}\dsep
  \bigcap M\neq\varnothing\vdash \Frankl{M}
}{%
  \DeltaRow{Mengen}{M\dsep X\dsep a}
  \DeltaRow{Funktionen}{F}
}
}

\expandafter\newcommand\csname FranklStatement@FranklPairAvoidADecomp\endcsname{%
\FormulaThmDeltaK[Zerlegung der \(a\)-vermeidenden Klassen im Zweierfall]{%
  \begin{aligned}[t]
  \FamAvoiding{M}{a}
  &=
  \bigl(\FamAvoiding{M}{a}\cap\FamAvoiding{M}{b}\bigr)
  \cup
  \bigl(\FamAvoiding{M}{a}\cap\FamContaining{M}{b}\bigr)
  \end{aligned}
}{FranklPairAvoidADecomp}{%
  \DeltaRow{Mengen}{M\dsep X\dsep a\dsep b}
}
}

\expandafter\newcommand\csname FranklStatement@FranklPairAvoidADisjoint\endcsname{%
\FormulaThmDeltaK[Disjunktheit der \(a\)-vermeidenden Zweige]{%
  \bigl(\FamAvoiding{M}{a}\cap\FamAvoiding{M}{b}\bigr)
  \cap
  \bigl(\FamAvoiding{M}{a}\cap\FamContaining{M}{b}\bigr)
  =\varnothing
}{FranklPairAvoidADisjoint}{%
  \DeltaRow{Mengen}{M\dsep X\dsep a\dsep b}
}
}

\expandafter\newcommand\csname FranklStatement@FranklPairContainADecomp\endcsname{%
\FormulaThmDeltaK[Zerlegung der \(a\)-enthaltenden Zielklassen]{%
  \begin{aligned}[t]
  \FamContaining{M}{a}
  &=
  \bigl(\FamContaining{M}{a}\cap\FamContaining{M}{b}\bigr)
  \cup
  \bigl(\FamContaining{M}{a}\cap\FamAvoiding{M}{b}\bigr)
  \end{aligned}
}{FranklPairContainADecomp}{%
  \DeltaRow{Mengen}{M\dsep X\dsep a\dsep b}
}
}

\expandafter\newcommand\csname FranklStatement@FranklPairContainADisjoint\endcsname{%
\FormulaThmDeltaK[Disjunktheit der \(a\)-enthaltenden Zielzweige]{%
  \bigl(\FamContaining{M}{a}\cap\FamContaining{M}{b}\bigr)
  \cap
  \bigl(\FamContaining{M}{a}\cap\FamAvoiding{M}{b}\bigr)
  =\varnothing
}{FranklPairContainADisjoint}{%
  \DeltaRow{Mengen}{M\dsep X\dsep a\dsep b}
}
}

\expandafter\newcommand\csname FranklStatement@FranklPairMixed01Finite\endcsname{%
\FormulaThmDeltaK[Endlichkeit der Klasse \(M_{01}\)]{%
  \Finite{M}\vdash
  \Finite{\FamAvoiding{M}{a}\cap\FamContaining{M}{b}}
}{FranklPairMixed01Finite}{%
  \DeltaRow{Mengen}{M\dsep a\dsep b}
}
}

\expandafter\newcommand\csname FranklStatement@FranklPairMixedReverseInjection\endcsname{%
\FormulaThmDeltaK[Gegeninjektion im negativen Mischklassenfall]{%
  \begin{aligned}[t]
  &\Finite{M}\dsep
  \neg\exists F\colon
  \bigl(\FamAvoiding{M}{a}\cap\FamContaining{M}{b}\bigr)
  \inj
  \bigl(\FamContaining{M}{a}\cap\FamAvoiding{M}{b}\bigr)
  \vdash{}\\
  &\exists G\colon
  \bigl(\FamContaining{M}{a}\cap\FamAvoiding{M}{b}\bigr)
  \inj
  \bigl(\FamAvoiding{M}{a}\cap\FamContaining{M}{b}\bigr)
  \end{aligned}
}{FranklPairMixedReverseInjection}{%
  \DeltaRow{Mengen}{M\dsep a\dsep b}
  \DeltaRow{Funktionen}{F\dsep G}
}
}

\expandafter\newcommand\csname FranklStatement@FranklPairAdjDef\endcsname{%
\FormulaDefDeltaK[Zweieradjunktionsabbildung]{%
  \UnionClosedFam(M)\dsep \{a,b\}\in M\vdash
  \begin{aligned}[t]
    &U_{M,a,b}\colon
      \bigl(\FamAvoiding{M}{a}\cap\FamAvoiding{M}{b}\bigr)
      \\[-2pt]
    &\hspace{4em}\to
      \bigl(\FamContaining{M}{a}\cap\FamContaining{M}{b}\bigr),\\[-2pt]
    &\forall X\in\FamAvoiding{M}{a}\cap\FamAvoiding{M}{b}\,
      \bigl(U_{M,a,b}(X)\coloneqq X\cup\{a,b\}\bigr)
  \end{aligned}
}{FranklPairAdjDef}{%
  \DeltaRow{Mengen}{M\dsep X\dsep a\dsep b}
  \DeltaRow{Funktionen}{U_{M,a,b}}
  \DeltaRow{\textbf{Neues Symbol}}{}
  \DeltaPrem{Zweieradjunktionsabbildung}{U_{M,a,b}}
}
}

\expandafter\newcommand\csname FranklStatement@FranklPairAdjFunction\endcsname{%
\FormulaThmDeltaK[Zweieradjunktionsabbildung ist Funktion]{%
  \UnionClosedFam(M)\dsep \{a,b\}\in M\vdash
  U_{M,a,b}\colon
  \bigl(\FamAvoiding{M}{a}\cap\FamAvoiding{M}{b}\bigr)
  \to
  \bigl(\FamContaining{M}{a}\cap\FamContaining{M}{b}\bigr)
}{FranklPairAdjFunction}{%
  \DeltaRow{Mengen}{M\dsep X\dsep a\dsep b}
  \DeltaRow{Funktionen}{U_{M,a,b}}
}
}

\expandafter\newcommand\csname FranklStatement@FranklPairAdjEval\endcsname{%
\FormulaThmDeltaK[Auswertung der Zweieradjunktionsabbildung]{%
  \UnionClosedFam(M)\dsep \{a,b\}\in M\dsep
  X\in \FamAvoiding{M}{a}\cap\FamAvoiding{M}{b}
  \vdash U_{M,a,b}(X)=X\cup\{a,b\}
}{FranklPairAdjEval}{%
  \DeltaRow{Mengen}{M\dsep X\dsep a\dsep b}
  \DeltaRow{Funktionen}{U_{M,a,b}}
}
}

\expandafter\newcommand\csname FranklStatement@FranklPairAdjInjection\endcsname{%
\FormulaThmDeltaK[Zweieradjunktionsabbildung ist injektiv]{%
  \UnionClosedFam(M)\dsep \{a,b\}\in M
  \vdash
  U_{M,a,b}\colon
  \bigl(\FamAvoiding{M}{a}\cap\FamAvoiding{M}{b}\bigr)
  \inj
  \bigl(\FamContaining{M}{a}\cap\FamContaining{M}{b}\bigr)
}{FranklPairAdjInjection}{%
  \DeltaRow{Mengen}{M\dsep X\dsep Y\dsep a\dsep b}
  \DeltaRow{Funktionen}{U_{M,a,b}}
}
}

\expandafter\newcommand\csname FranklStatement@FranklPairPositiveBranch\endcsname{%
\FormulaThmDeltaK[Positiver Mischklassenfall liefert Frankl]{%
  \begin{aligned}[t]
  &\UnionClosedFam(M)\dsep \{a,b\}\in M\dsep a\in\bigcup M\dsep{}\\
  &\exists F\colon
  \bigl(\FamAvoiding{M}{a}\cap\FamContaining{M}{b}\bigr)
  \inj
  \bigl(\FamContaining{M}{a}\cap\FamAvoiding{M}{b}\bigr)
  \vdash \Frankl{M}
  \end{aligned}
}{FranklPairPositiveBranch}{%
  \DeltaRow{Mengen}{M\dsep a\dsep b}
  \DeltaRow{Funktionen}{F\dsep H\dsep U_{M,a,b}}
}
}

\expandafter\newcommand\csname FranklStatement@FranklPairNegativeBranch\endcsname{%
\FormulaThmDeltaK[Negativer Mischklassenfall liefert Frankl]{%
  \begin{aligned}[t]
  &\Finite{M}\dsep \UnionClosedFam(M)\dsep
  \{a,b\}\in M\dsep b\in\bigcup M\dsep{}\\
  &\neg\exists F\colon
  \bigl(\FamAvoiding{M}{a}\cap\FamContaining{M}{b}\bigr)
  \inj
  \bigl(\FamContaining{M}{a}\cap\FamAvoiding{M}{b}\bigr)
  \vdash \Frankl{M}
  \end{aligned}
}{FranklPairNegativeBranch}{%
  \DeltaRow{Mengen}{M\dsep a\dsep b}
  \DeltaRow{Funktionen}{F\dsep G}
}
}

\expandafter\newcommand\csname FranklStatement@FranklPairMemberFrankl\endcsname{%
\FormulaThmDeltaK[Frankl-Vermutung bei enthaltener Zweiermenge]{%
  \Finite{M}\dsep \UnionClosedFam(M)\dsep a\neq b\dsep \{a,b\}\in M
  \vdash \Frankl{M}
}{FranklPairMemberFrankl}{%
  \DeltaRow{Mengen}{M\dsep a\dsep b}
  \DeltaRow{Funktionen}{F}
}
}

\expandafter\newcommand\csname FranklStatement@FranklFamPairMemberFrankl\endcsname{%
\FormulaThmDeltaK[Frankl-Familie mit enthaltener Zweiermenge erfüllt Frankl]{%
  \FranklFam{M}\dsep
  \exists a\,\exists b\,(a\neq b\land \{a,b\}\in M)
  \vdash \Frankl{M}
}{FranklFamPairMemberFrankl}{%
  \DeltaRow{Mengen}{M\dsep a\dsep b}
}
}

